\documentclass{article}
\usepackage[T1]{fontenc}
\usepackage{lmodern}
\usepackage{graphicx}
\usepackage{amsmath,amssymb}
\usepackage[a4paper,margin=2cm]{geometry}
\usepackage[absolute,overlay]{textpos}
\setlength{\TPHorizModule}{1mm}
\setlength{\TPVertModule}{1mm}
\usepackage[indonesian]{babel}
\usepackage{hyperref}
\setlength{\emergencystretch}{2em}
% unit: Halaman 10

\begin{document}

% === Halaman 10 (python/native) ===

4) Bobot total sekarang menjadi 18 – 13 = 5.

5) Bandingkan 5 dengan bobot terbesar kedua, yaitu 6. Karena 6 > 5, 

maka 6 tidak dimasukkan ke dalam knapsack. → b3 = 0

6) Bandingkan 5 dengan bobot terbesar berikutnya, yaitu 3.  Karena 

3  5, maka 3 dimasukkan ke dalam knapsack. → b2 = 1

7) Bobot total sekarang menjadi 5 – 3 = 2.

8) Bandingkan 2 dengan bobot terbesar berikutnya, yaitu 2. Karena 2 

 2, maka 2 dimasukkan ke dalam knapsack. → b1 = 0

9) Bobot total sekarang menjadi 2 – 2 = 0.

10

\end{document}
